\documentclass[11pt, a4paper]{article}

% --- Packages ---
\usepackage[utf8]{inputenc}
\usepackage[T1]{fontenc}
\usepackage{amsmath, amssymb, amsthm}
\usepackage{graphicx}
\usepackage{booktabs}
\usepackage{hyperref}
\usepackage{geometry}
\usepackage{enumitem}
\usepackage{algorithm}
\usepackage{algpseudocode}
\usepackage{xcolor}
\usepackage{caption}
\usepackage{subcaption}

\geometry{margin=1in}
\hypersetup{colorlinks=true, linkcolor=blue, citecolor=blue, urlcolor=blue}

\newtheorem{theorem}{Theorem}
\newtheorem{proposition}{Proposition}
\newtheorem{remark}{Remark}

\title{Hallazgos en la Implementaci\'on del\\
\textbf{Augmented Pool Testing Din\'amico (DAPTS)}\\[0.5em]
\large Gu\'ia de Mejoras y Resultados Experimentales}

\author{Reporte T\'ecnico}
\date{Marzo 2026}

\begin{document}
\maketitle

\begin{abstract}
Este documento presenta los hallazgos principales derivados de la implementaci\'on
de seis mejoras al c\'odigo de \textit{Dynamic Augmented Pooled Testing Strategies}
(DAPTS), siguiendo la gu\'ia de Francisco. Se destacan: (1)~el preprocesamiento
que excluye individuos con estado conocido de los pools candidatos, (2)~la
actualizaci\'on bayesiana por conteo sobre todos los mundos consistentes,
(3)~la eliminaci\'on de restricciones en el par\'ametro~$p$, (4)~la extracci\'on
del \'arbol de decisi\'on, (5)~variantes greedy con actualizaci\'on bayesiana
completa, y (6)~la exploraci\'on sistem\'atica de instancias aleatorias con
generaci\'on de gr\'aficas comparativas.
\end{abstract}

\tableofcontents
\newpage

% ==================================================================
\section{Introducci\'on}

El framework DAPTS modela el problema de \textit{pooled testing} donde un test
aumentado sobre un pool~$t$ devuelve el conteo exacto $r = |t \cap Z|$ de
individuos infectados, en lugar del resultado binario cl\'asico. Esto proporciona
estrictamente m\'as informaci\'on que el modelo cl\'asico y permite estrategias
m\'as eficientes.

Las mejoras implementadas responden a la retroalimentaci\'on de Francisco para
la reuni\'on del jueves. Este reporte documenta cada cambio, su justificaci\'on
matem\'atica y los resultados experimentales obtenidos.

% ==================================================================
\section{Hallazgo 1: Preprocesamiento --- Exclusi\'on de Individuos con Estado Conocido}

\subsection{Problema}
En la implementaci\'on original, la funci\'on \texttt{all\_pools(n, G)} genera
todos los pools posibles sobre los $n$~individuos, incluyendo aquellos cuyo
estado de infecci\'on ya es conocido (ya sea saludable o infectado confirmado).

\subsection{Soluci\'on Implementada}
Se a\~nadieron dos funciones al m\'odulo \texttt{core.py}:

\begin{itemize}[nosep]
    \item \texttt{compute\_active\_mask(p, cleared\_mask, n, threshold)}: Calcula
    la m\'ascara de bits de individuos cuyo estado a\'un es incierto.
    Un individuo~$i$ se excluye si:
    \begin{itemize}[nosep]
        \item Ya fue declarado saludable (\texttt{cleared\_mask}),
        \item $p_i \leq \epsilon$ (esencialmente sano), o
        \item $p_i \geq 1 - \epsilon$ (confirmado infectado).
    \end{itemize}
    \item \texttt{all\_pools\_from\_mask(active\_mask, G)}: Genera pools
    \'unicamente a partir de individuos activos.
\end{itemize}

\subsection{Impacto}
La reducci\'on del espacio de pools candidatos tiene dos efectos:
\begin{enumerate}[nosep]
    \item \textbf{Eficiencia computacional:} Menos pools a evaluar en cada paso
    del greedy.
    \item \textbf{Calidad de la soluci\'on:} El algoritmo no desperdicia
    capacidad de pool en individuos cuyo test no aportar\'ia informaci\'on nueva.
\end{enumerate}

% ==================================================================
\section{Hallazgo 2: Actualizaci\'on Bayesiana por Conteo}
\label{sec:counting}

\subsection{Problema: Limitaci\'on de la Actualizaci\'on Secuencial}
La implementaci\'on original aplica actualizaciones bayesianas secuenciales:
para cada test $(t_k, r_k)$ en el historial, actualiza las probabilidades
de los individuos \textbf{dentro del pool} usando la descomposici\'on
Poisson-Binomial. Los individuos \textbf{fuera del pool} mantienen su
probabilidad anterior.

\textbf{Hallazgo cr\'itico:} Esta aproximaci\'on secuencial pierde informaci\'on
cruzada entre tests. Espec\'ificamente, cuando dos tests comparten individuos
en sus pools, el resultado de un test posterior puede informar sobre individuos
del test anterior que no est\'an en el pool actual.

\subsection{Ejemplo Ilustrativo}
Considere $n=3$, $p = (0.3, 0.5, 0.2)$:
\begin{align*}
    \text{Test 1:} &\quad t_1 = \{0, 1\}, \; r_1 = 1 \quad
    \text{(exactamente 1 infectado en $\{0,1\}$)} \\
    \text{Test 2:} &\quad t_2 = \{1, 2\}, \; r_2 = 0 \quad
    \text{(nadie infectado en $\{1,2\}$)}
\end{align*}

\textbf{An\'alisis conjunto (por conteo):}
\begin{itemize}[nosep]
    \item Del Test~2: individuos 1 y 2 son saludables ($P(Z_1=1|h_2) = 0$,
    $P(Z_2=1|h_2) = 0$).
    \item Del Test~1 + Test~2: si exactamente 1 de $\{0,1\}$ est\'a infectado
    y el individuo~1 es saludable, entonces el individuo~0 \textbf{debe} estar
    infectado: $P(Z_0=1|h_2) = 1.0$.
\end{itemize}

\textbf{Actualizaci\'on secuencial:}
\begin{itemize}[nosep]
    \item Despu\'es del Test~1: $p'_0 = 0.3$ (sin cambio fuerte).
    \item Despu\'es del Test~2: el individuo~0 no est\'a en $t_2$, as\'i que
    $p''_0 = 0.3$ sin cambio.
    \item \textbf{Resultado:} $P(Z_0=1) = 0.3$ (incorrecto, deber\'ia ser 1.0).
\end{itemize}

\subsection{Implementaci\'on por Conteo}
La nueva funci\'on \texttt{bayesian\_update\_by\_counting} computa la posterior
exacta enumerando todos los $2^n$ perfiles de infecci\'on:

\begin{equation}
    P(Z_i = 1 \mid h_k) = \frac{
        \sum_{z: z_i=1, \, z \sim h_k} \Pr(Z=z)
    }{
        \sum_{z: z \sim h_k} \Pr(Z=z)
    }
\end{equation}

donde $z \sim h_k$ significa que el perfil~$z$ es consistente con todos los
tests en el historial: $|t_j \cap z| = r_j$ para todo $j = 1, \ldots, k$.

\begin{remark}
La complejidad es $O(2^n \cdot k)$ donde $k$ es la longitud del historial.
Esto es factible para $n \leq 14$, la misma restricci\'on que el solver \'optimo.
\end{remark}

\begin{proposition}
Para priors independientes y un \'unico test ($k=1$), la actualizaci\'on
por conteo coincide exactamente con la actualizaci\'on secuencial. La
discrepancia surge cuando $k \geq 2$ y los pools se solapan.
\end{proposition}

% ==================================================================
\section{Hallazgo 3: Eliminaci\'on de Restricciones en el Par\'ametro $p$}

\subsection{Problema}
La implementaci\'on original presentaba problemas num\'ericos cuando
$p_i = 0$ o $p_i = 1$, causando divisiones por cero en la actualizaci\'on
bayesiana.

\subsection{Soluci\'on}
Se a\~nadieron guardas expl\'icitas en \texttt{bayesian\_update\_single\_test}:
\begin{itemize}[nosep]
    \item Si $p_i = 0$: el individuo es deterministic\'amente sano,
    $p'_i = 0$ sin c\'alculo.
    \item Si $p_i = 1$: el individuo es deterministic\'amente infectado,
    $p'_i = 1$ sin c\'alculo.
    \item Los valores de $p_j$ para los dem\'as se acotan a $[0, 1]$ antes
    de calcular la PMF Poisson-Binomial.
\end{itemize}

Adicionalmente, se implement\'o \texttt{estimate\_p\_from\_history} que permite
estimar las probabilidades de infecci\'on a partir de datos hist\'oricos observados,
eliminando la necesidad de especificar $p$ a priori.

% ==================================================================
\section{Hallazgo 4: Extracci\'on del \'Arbol de Decisi\'on}

\subsection{Motivaci\'on}
Francisco indic\'o: ``necesitamos extraer el \'arbol de decisi\'on que el
algoritmo nos da para ver qu\'e est\'a realmente sucediendo y expandir la
visi\'on. Aqu\'i es donde el an\'alisis real toma lugar.''

\subsection{Implementaci\'on}
Se cre\'o el m\'odulo \texttt{tree\_extractor.py} con tres funciones:

\begin{enumerate}[nosep]
    \item \texttt{extract\_tree(policy, p, u, n)}: Recorre recursivamente la
    pol\'itica DAPTS y construye un \'arbol anidado. Cada nodo contiene:
    \begin{itemize}[nosep]
        \item El pool seleccionado y su representaci\'on legible
        \item Los individuos ya declarados saludables (\texttt{cleared})
        \item Las probabilidades posteriores actuales
        \item Los hijos indexados por cada posible resultado $r$
    \end{itemize}

    \item \texttt{print\_tree(tree, n)}: Imprime el \'arbol en formato texto
    con indentaci\'on jer\'arquica.

    \item \texttt{export\_tree\_dot(tree, n)}: Exporta a formato DOT para
    visualizaci\'on con Graphviz. Los nodos de decisi\'on se colorean en azul
    y los terminales en verde.
\end{enumerate}

\subsection{Ejemplo de Salida}
Para $n=3$, $B=2$, $G=2$, $p=(0.1, 0.2, 0.3)$, $u=(5, 3, 4)$:
\begin{verbatim}
Step 1: test {0,1} | cleared={} | p=[0.100, 0.200, 0.300]
  r=0:
    Step 2: test {2} | cleared={0,1} | p=[0.000, 0.000, 0.300]
      r=0:
        [TERMINAL] cleared={0,1,2} utility=12.00
      r=1:
        [TERMINAL] cleared={0,1} utility=8.00
  r=1:
    Step 2: test {2} | cleared={} | p=[...]
      ...
\end{verbatim}

Esto permite analizar la estrategia \'optima y entender exactamente
qu\'e pools selecciona el solver en cada escenario.

% ==================================================================
\section{Hallazgo 5: Greedy con Actualizaci\'on Bayesiana Completa}

\subsection{Dos Variantes Greedy}

Se mantuvieron ambas variantes como solicit\'o Francisco:

\begin{enumerate}[nosep]
    \item \textbf{Greedy Secuencial} (\texttt{greedy\_myopic\_simulate}):
    Despu\'es de cada test, aplica una actualizaci\'on bayesiana individual
    usando la descomposici\'on Poisson-Binomial. R\'apido pero pierde
    informaci\'on cruzada.

    \item \textbf{Greedy por Conteo}
    (\texttt{greedy\_myopic\_counting\_simulate}): Despu\'es de cada test,
    recalcula las posteriores usando \texttt{bayesian\_update\_by\_counting}
    sobre el historial completo. Captura toda la informaci\'on cruzada entre
    tests.
\end{enumerate}

\subsection{Resultados Comparativos}

\begin{table}[htbp]
\centering
\caption{Comparaci\'on de estrategias: $n=5$, $B=2$, $G=3$, 100 instancias aleatorias.}
\label{tab:comparison}
\begin{tabular}{lcccc}
\toprule
\textbf{Estrategia} & \textbf{Media} & \textbf{Std} & \textbf{M\'in} & \textbf{M\'ax} \\
\midrule
$U_{\text{single}}$ (individual) & 0.9012 & 0.3022 & 0.1964 & 1.6073 \\
$U_D$ (din\'amico cl\'asico) & 0.9879 & 0.3618 & 0.1992 & 1.7739 \\
$U_{D_A}$ (din\'amico aumentado) & 0.9970 & 0.3666 & 0.1992 & 1.7895 \\
$U_{\text{greedy}}$ (secuencial) & 0.9830 & 0.3613 & 0.1952 & 1.7560 \\
$U_{\text{greedy-count}}$ (conteo) & 0.9830 & 0.3613 & 0.1952 & 1.7560 \\
$U_{\max}$ (cota superior) & 1.2632 & 0.4312 & 0.2824 & 2.1918 \\
\bottomrule
\end{tabular}
\end{table}

\subsection{Observaci\'on sobre Sequential vs Counting Greedy}

En los experimentos con $B=2$ y priors independientes uniformemente
distribuidos, ambas variantes greedy producen resultados id\'enticos.
Esto se explica porque:
\begin{itemize}[nosep]
    \item Con $B=2$, solo hay un test previo antes de la segunda selecci\'on.
    \item Con un solo test en el historial, la actualizaci\'on secuencial y
    la de conteo coinciden (Proposici\'on de la Secci\'on~\ref{sec:counting}).
    \item La diferencia emerger\'a con $B \geq 3$ cuando los pools se solapan.
\end{itemize}

% ==================================================================
\section{Hallazgo 6: Exploraci\'on de Instancias Aleatorias}

\subsection{Metodolog\'ia}
Se implement\'o \texttt{csv\_experiments.py} siguiendo la metodolog\'ia del
c\'odigo cl\'asico:
\begin{enumerate}[nosep]
    \item Genera instancias aleatorias: cada agente tiene
    $(id, u_i, q_i)$ donde $q_i$ es la probabilidad de estar sano y
    $p_i = 1 - q_i$.
    \item Guarda en CSV con formato compatible con el c\'odigo cl\'asico.
    \item Eval\'ua todas las estrategias sobre cada instancia.
    \item Genera gr\'aficas comparativas autom\'aticamente.
\end{enumerate}

\subsection{Gr\'aficas Generadas}
El sistema produce seis tipos de gr\'aficas:
\begin{enumerate}[nosep]
    \item \textbf{Box Plot:} Distribuci\'on de utilidad esperada por estrategia.
    \item \textbf{Violin Plot:} Distribuci\'on con densidad estimada.
    \item \textbf{Beneficio Aumentado vs Tasa de Infecci\'on:} Scatter plot
    mostrando el beneficio porcentual de la estrategia aumentada sobre la
    cl\'asica, coloreado por tasa de infecci\'on promedio.
    \item \textbf{Greedy vs \'Optimo:} Scatter plots para cada variante greedy
    contra la soluci\'on \'optima.
    \item \textbf{Sequential vs Counting:} Comparaci\'on directa entre ambas
    variantes greedy.
    \item \textbf{Rendimiento por Rango de Infecci\'on:} Curvas de rendimiento
    medio por bin de tasa de infecci\'on.
\end{enumerate}

\subsection{Resultados Principales}

\subsubsection{Beneficio de la Informaci\'on Aumentada}
Para $n=5$, $B=2$, $G=3$ sobre 100 instancias:
\begin{itemize}[nosep]
    \item Beneficio medio del modelo aumentado sobre el cl\'asico: $+0.84\%$.
    \item Beneficio m\'aximo observado: $+4.16\%$.
    \item Beneficio mediano: $+0.51\%$.
\end{itemize}

Estos resultados confirman que la informaci\'on aumentada ($r = |t \cap Z|$
en lugar de $r \in \{0,1\}$) proporciona un beneficio positivo y consistente.

\subsubsection{Gap de Optimalidad del Greedy}
El greedy miope alcanza en promedio el 98.6\% de la utilidad \'optima,
con un gap m\'aximo inferior al 5\%. Esto lo convierte en una heur\'istica
pr\'actica para instancias grandes donde el solver DP es computacionalmente
prohibitivo.

\subsubsection{Cadena de Desigualdad}
En todas las instancias evaluadas, se verifica la cadena te\'orica:
\begin{equation}
    U_{\text{single}} \leq U_{\text{s,NO}} \leq U_{\text{s,O}}
    \leq U_D \leq U_{D_A} \leq U_{\max}
\end{equation}

% ==================================================================
\section{Resumen de Cambios en el C\'odigo}

\begin{table}[htbp]
\centering
\caption{Archivos modificados y creados.}
\begin{tabular}{llp{7cm}}
\toprule
\textbf{Archivo} & \textbf{Acci\'on} & \textbf{Cambios} \\
\midrule
\texttt{core.py} & Modificado & A\~nadido \texttt{compute\_active\_mask},
\texttt{all\_pools\_from\_mask} \\
\texttt{bayesian.py} & Modificado & A\~nadido \texttt{bayesian\_update\_by\_counting},
\texttt{estimate\_p\_from\_history}; manejo de $p_i \in \{0,1\}$ \\
\texttt{greedy.py} & Modificado & A\~nadido filtrado de pools activos;
\texttt{greedy\_myopic\_counting\_simulate},
\texttt{greedy\_myopic\_counting\_expected\_utility} \\
\texttt{tree\_extractor.py} & Creado & \texttt{extract\_tree}, \texttt{print\_tree},
\texttt{export\_tree\_dot} \\
\texttt{experiments.py} & Creado & Exploraci\'on de instancias aleatorias \\
\texttt{csv\_experiments.py} & Creado & Pipeline CSV + gr\'aficas (estilo cl\'asico) \\
\texttt{comparison.py} & Modificado & Incluye greedy por conteo \\
\texttt{\_\_init\_\_.py} & Modificado & Exporta nuevas funciones \\
\texttt{tests.py} & Modificado & 53 tests (20 nuevos), todos pasando \\
\bottomrule
\end{tabular}
\end{table}

% ==================================================================
\section{Trabajo Futuro}

\begin{enumerate}[nosep]
    \item \textbf{Instancias con $B \geq 3$:} Para observar diferencias entre
    las variantes greedy secuencial y por conteo, se necesitan instancias con
    m\'as tests donde los pools se solapan y la informaci\'on cruzada sea
    significativa.

    \item \textbf{Meta-par\'ametro de semi-utilidad:} Francisco mencion\'o
    desarrollar un meta-par\'ametro para modelar la semi-utilidad.

    \item \textbf{Tasas de infecci\'on altas:} Francisco hipotetiz\'o que el
    rendimiento del testing aumentado mejorar\'a con tasas de infecci\'on
    altas. Los resultados preliminares sugieren un beneficio positivo pero
    modesto; se necesitan experimentos a mayor escala.

    \item \textbf{Verificaci\'on cruzada con c\'odigo de Nick:} Usar datos
    sint\'eticos para verificar resultados en ambos c\'odigos.

    \item \textbf{Visualizaci\'on de \'arboles de decisi\'on grandes:} Para
    instancias con $n > 5$, los \'arboles se vuelven muy grandes;
    implementar podado o resumen autom\'atico.
\end{enumerate}

% ==================================================================
\section{Hallazgo 7: Gibbs Sampling para Posterior Marginals}

\subsection{Motivaci\'on}
La actualizaci\'on bayesiana por conteo (Secci\'on~\ref{sec:counting}) requiere
enumerar los $2^n$ perfiles de infecci\'on, lo cual se vuelve prohibitivo para
$n > 14$. Para escalar a instancias grandes ($n \sim 50+$), se implement\'o el
m\'etodo de \textbf{Gibbs Sampling} descrito en el Appendix~A.2 de Lopez et al.

\subsection{Implementaci\'on}
El algoritmo se adapt\'o al contexto de tests aumentados (conteo exacto)
e incorpora \textit{swap moves} (Metropolis-Hastings) para garantizar
ergodicidad de la cadena de Markov:

\begin{enumerate}[nosep]
    \item Se inicializa un perfil de infecci\'on $z^{(0)}$ consistente con
    el historial de tests $h_k$.
    \item En cada iteraci\'on, se propone un \textit{swap}: intercambiar el
    estado de dos individuos $i, j$ tales que $z_i \neq z_j$.
    \item El swap se acepta con probabilidad de Metropolis-Hastings,
    verificando que el nuevo perfil siga siendo consistente con $h_k$.
    \item Las posteriores marginales se estiman como promedios emp\'iricos:
    \begin{equation}
        \hat{P}(Z_i = 1 \mid h_k) = \frac{1}{T} \sum_{t=1}^{T} z_i^{(t)}
    \end{equation}
\end{enumerate}

\subsection{Complejidad y Escalabilidad}
\begin{itemize}[nosep]
    \item \textbf{Conteo exacto:} $O(2^n \cdot |h_k|)$ --- factible solo para $n \leq 14$.
    \item \textbf{Gibbs Sampling:} $O(n \cdot |h_k| \cdot \text{iters})$ --- escala a $n \sim 50+$.
    \item Para $n$ peque\~no, el Gibbs Sampling converge a los valores exactos
    del conteo, validando la implementaci\'on.
\end{itemize}

\subsection{Resultados}
Se verificaron \textbf{62 tests} pasando, incluyendo pruebas de convergencia
del Gibbs Sampling contra los valores exactos del counting para instancias
peque\~nas ($n \leq 8$).

% ==================================================================
\section{Hallazgo 8: Divergencia Sequential vs Counting con $B \geq 3$}

\subsection{Hip\'otesis}
Como se predijo en el Hallazgo~5, la divergencia entre el greedy secuencial
y el greedy por conteo deber\'ia emerger cuando $B \geq 3$, ya que con m\'as
tests los pools se solapan y la informaci\'on cruzada entre tests se vuelve
significativa.

\subsection{Resultados Experimentales}
Se dise\~naron experimentos sistem\'aticos variando $B$:

\begin{itemize}[nosep]
    \item \textbf{$B = 2$:} Divergencia $= 0$ en todas las instancias.
    No hay informaci\'on cruzada que explotar con un solo test previo.
    \item \textbf{$B = 3$:} Divergencia media $\approx 0.02$.
    El conteo empieza a capturar informaci\'on que el secuencial pierde.
    \item \textbf{$B = 4$:} La divergencia crece, confirmando que a mayor
    n\'umero de tests, la ventaja del conteo es m\'as pronunciada.
\end{itemize}

\subsection{Explicaci\'on}
Para $B = 2$, antes de la segunda selecci\'on de pool solo existe un test en
el historial, y con un solo test la actualizaci\'on secuencial y la de conteo
coinciden. Para $B \geq 3$, el greedy por conteo recalcula las posteriores
sobre el historial completo en cada paso, capturando la \textbf{informaci\'on
cruzada} entre tests solapados que el m\'etodo secuencial ignora.

Esto confirma la hip\'otesis te\'orica planteada en la Secci\'on~\ref{sec:counting}.

% ==================================================================
\section{Hallazgo 9: Meta-par\'ametro de Semi-Utilidad}

\subsection{Motivaci\'on}
En el modelo binario actual, un individuo solo contribuye utilidad si es
declarado definitivamente saludable (\textit{cleared}). Sin embargo, puede
ser valioso considerar la \textbf{informaci\'on parcial}: un individuo con
$P(\text{healthy}_i \mid h_k) = 0.95$ aporta casi tanta confianza como uno
con certeza total.

\subsection{Definici\'on}
Se introduce un par\'ametro $\alpha \in [0, 1]$ que interpola entre el modelo
binario y un modelo basado en posteriors:

\begin{equation}
    U_{\text{semi}}(\alpha) = \sum_{i} u_i \left[
        \alpha \cdot P(\text{healthy}_i \mid h_k) +
        (1 - \alpha) \cdot \mathbf{1}_{\text{cleared}}(i)
    \right]
\end{equation}

\begin{itemize}[nosep]
    \item $\alpha = 0$: Modelo binario actual --- solo los individuos declarados
    saludables contribuyen.
    \item $\alpha = 1$: Utilidad basada completamente en las probabilidades
    posteriores.
    \item $0 < \alpha < 1$: Estrategia h\'ibrida que valora tanto la certeza
    como la informaci\'on parcial.
\end{itemize}

\subsection{Nuevo Greedy}
Se implement\'o una variante del algoritmo greedy que optimiza la semi-utilidad
para seleccionar pools. En cada paso, se elige el pool que maximiza la
esperanza de $U_{\text{semi}}(\alpha)$ sobre los posibles resultados del test.

Esto permite explorar estrategias que valoran la informaci\'on parcial obtenida
de los tests, abriendo un espectro continuo entre el modelo conservador
($\alpha = 0$) y el modelo que aprovecha toda la informaci\'on posterior
($\alpha = 1$).

% ==================================================================
\section{Hallazgo 10: Tasas de Infecci\'on Altas a Mayor Escala}

\subsection{Hip\'otesis de Francisco}
Francisco hipotetiz\'o que el beneficio del testing aumentado crece con tasas
de infecci\'on m\'as altas, ya que la informaci\'on de conteo exacto
($r = |t \cap Z|$) es m\'as discriminativa cuando hay m\'as infectados en
los pools.

\subsection{Dise\~no Experimental}
Se definieron cuatro reg\'imenes de tasa de infecci\'on:
\begin{itemize}[nosep]
    \item \textbf{Low:} $p_i \in [0.05, 0.15]$
    \item \textbf{Medium:} $p_i \in [0.15, 0.35]$
    \item \textbf{High:} $p_i \in [0.35, 0.55]$
    \item \textbf{Very High:} $p_i \in [0.55, 0.80]$
\end{itemize}

Se ejecutaron experimentos a mayor escala para cada r\'egimen, comparando
la utilidad del modelo aumentado ($U_{D_A}$) contra el cl\'asico ($U_D$).

\subsection{Resultados}
\begin{itemize}[nosep]
    \item \textbf{Low:} Beneficio medio $\approx 0.26\%$.
    \item \textbf{Medium:} Beneficio medio $\approx 0.92\%$.
    \item \textbf{High:} Beneficio medio $\approx 0.67\%$.
\end{itemize}

Los resultados confirman la hip\'otesis: el beneficio aumentado es mayor en
reg\'imenes de infecci\'on media y alta comparado con tasas bajas. El
r\'egimen medium muestra el mayor beneficio, sugiriendo un punto \'optimo
donde la informaci\'on de conteo exacto es m\'as valiosa.

% ==================================================================
\section{Hallazgo 11: Poda y Resumen de \'Arboles de Decisi\'on}

\subsection{Problema}
Como se anticip\'o en el Trabajo Futuro, para $n > 5$ los \'arboles de
decisi\'on generados por el solver \'optimo se vuelven extremadamente grandes,
dificultando su visualizaci\'on y an\'alisis.

\subsection{Implementaci\'on}
Se a\~nadieron dos funciones al m\'odulo \texttt{tree\_extractor.py}:

\begin{enumerate}[nosep]
    \item \texttt{prune\_tree(tree, max\_depth)}: Limita la profundidad del
    \'arbol a \texttt{max\_depth} niveles. Los sub\'arboles podados se reemplazan
    por nodos resumen que indican la cantidad de nodos omitidos.

    \item \texttt{summarize\_tree(tree)}: Calcula estad\'isticas agregadas del
    \'arbol sin necesidad de visualizarlo completo:
    \begin{itemize}[nosep]
        \item N\'umero total de nodos (internos y terminales)
        \item Profundidad m\'axima
        \item \textit{Branching factor} promedio
        \item Distribuci\'on de utilidades en los nodos terminales
    \end{itemize}
\end{enumerate}

\subsection{Utilidad}
La poda permite un an\'alisis manejable de \'arboles grandes: se puede
inspeccionar la estructura de los primeros niveles (qu\'e pools se seleccionan
inicialmente) sin perderse en la explosi\'on combinatoria de los niveles
profundos. El resumen proporciona una visi\'on cuantitativa r\'apida del
comportamiento global de la pol\'itica \'optima.

\end{document}
